\section{First UIPs}

Our next addition to CDCL is First UIPs. UIP stands for "unit implication point," which 
is called a ``dominator'' in our implication graph. What is a dominator? A node $d$ is said 
to \textit{dominate} another node $n$ if every path in the graph from the entry node to $n$ 
passes through $d$. Think of a binary tree. Our entry node is the the head of the tree, and 
it dominates its two children. Likewise, each of those children dominates their own children 
nodes. So in other words, the UIP in our implication graph is kind of like the reason why 
we had a contradiction in the first place!

Given that summary of what a UIP is, wouldn't the first UIP be the decision we made on 
the current decision level? After all, we wouldn't have had a contradiction if we didn't 
make that decision. Well not necessarily, because in our implication graph, every decision 
we make only has outward edges, whereas all propagated nodes have inward edges as well. 
So while it is true that a first UIP could be a decision variable, it's not true that it's 
\textit{always} a decision variable. Take this previous implication graph as an example:

\noindent
\fbox{
\resizebox{\textwidth}{!}{%
\begin{tikzpicture}[
    >=Stealth,
    every node/.style={font=\large},
    implication/.style={->, thick}
]

% Left side
\node (x31) {$x_{31}=0\ @\ 3$};

\node (x2)
    [below right=0.8cm and 1.8cm of x31]
    {$x_2=0\ @\ 5$};

\node (x1)
    [below left=0.8cm and 1.8cm of x2]
    {$x_1=0\ @\ 5$};

\node (x3)
    [below right=0.8cm and 1.8cm of x1]
    {$x_3=0\ @\ 5$};

\node (x4)
    [above right=0.8cm and 1.8cm of x3]
    {$x_4=1\ @\ 5$};

% Right side
\node (x5)
    [above right=0.8cm and 2.5cm of x4]
    {$x_5=0\ @\ 5$};

\node (x6)
    [below right=0.8cm and 2.5cm of x4]
    {$x_6=0\ @\ 5$};

\node (x21)
    [below left=1.0cm and 1.8cm of x6]
    {$x_{21}=0\ @\ 2$};

\node (kappa)
    [right=2.5cm of x6]
    {$\kappa$};

% Left implication graph
\draw[implication]
    (x31) -- node[above left] {$\omega_1$} (x2);

\draw[implication]
    (x1) -- node[above left] {$\omega_1$} (x2);

\draw[implication]
    (x1) -- node[below left] {$\omega_2$} (x3);

\draw[implication]
    (x2) -- node[above right] {$\omega_3$} (x4);

\draw[implication]
    (x3) -- node[below right] {$\omega_3$} (x4);

% Right implication graph
\draw[implication]
    (x4) -- node[above left] {$\omega_4$} (x5);

\draw[implication]
    (x4) -- node[below left] {$\omega_5$} (x6);

\draw[implication]
    (x21) -- node[below right] {$\omega_5$} (x6);

\draw[implication]
    (x5) -- node[above right] {$\omega_6$} (kappa);

\draw[implication]
    (x6) -- node[below right] {$\omega_6$} (kappa);

\end{tikzpicture}
}
}
$x_4$ is a UIP in our implication graph because every path to $x_5$ and $x_6$ all 
pass through $x_4$'s node.

So what does the first UIP offer us? Resolving clauses up until we encounter the first UIP 
maximizes the level we backjump to and it makes the clauses we learn shorter, not just in length,
but also in how long it takes to resolve a clause. More details can be found here:
\begin{center}
    \small{\url{https://www3.cs.stonybrook.edu/~cram/cse505/Fall20/Resources/cdcl.pdf}}
\end{center}

How do we find the first UIP? The first UIP in a conflict is found when a the number 
of literals in the clause for the \textit{current} decision level is 1. So if our current 
decision level is 6, we resolve only on literals from decision level 6 and stop once the 
clause contains only 1 literal from decision level 6.

\subsection{Implementing First UIPs}

Once again, we only need to touch the conflict resolution code:

\begin{algorithm}[H]
\caption{Naive First UIP Conflict Resolution}\label{alg:33}
\begin{algorithmic}[1]
\Procedure{resolveConflict}{conflict[1..n]}

\State curLevel $\gets$ decisionLevel
\State end $\gets$ trail size
\State begin $\gets$ 0
\State resolveOn $\gets$ 0

\For{i $\gets 1$ up to $n$}
    \State var $\gets$ conflict[i]
    \State idx $\gets \Call{abs}{var}$
    \State tIdx $\gets$ varToTrail[idx]
    \If{trail[tIdx].decisionLevel = curLevel}
        \State end $\gets \Call{min}{end, tIdx}$
        \State begin $\gets \Call{max}{begin, tIdx}$
        \If{trail[tIdx].reasonIdx exists}
            \State resolveOn $\gets$ trail[tIdx].lit
        \EndIf
    \EndIf
\EndFor

\While{begin $\neq$ end}
    \State lit $\gets \Call{abs}{resolveOn}$
    \State tIdx $\gets$ varToTrail[resolveOn]
    \State reason $\gets$ Clauses[trail[tIdx].reasonIdx]
    \State $\Call{append}{conflict, reason}$

    \State remove all instances of lit and -lit from $conflict$
    \State remove all duplicate literals from $conflict$

    \For{each $var$ in conflict}
        \State tIdx $\gets$ varToTrail[$\Call{abs}{var}$]
        \If{trail[tIdx].decisionLevel = curLevel}
            \State end $\gets \Call{min}{end, tIdx}$
            \State begin $\gets \Call{max}{begin, tIdx}$
            \If{trail[tIdx].reasonIdx exists}
                \State resolveOn $\gets$ trail[tIdx].lit
            \EndIf
        \EndIf
    \EndFor
\EndWhile

\EndProcedure
\end{algorithmic}
\end{algorithm}

This is not a very efficient way of computing the first UIP (which you'll see in the 
results). However, let's analyze it anyways. It's the normal resolution code, but we 
do something different, which is stop when there's only one literal from the current 
decision level left in the clause. In our pseudocode, we don't compute the counts of 
literals from the current decision level, we have the condition $begin != end$. Why is 
that equivalent? Well \textit{end} is the index of the least recently propagated literal 
in the resolution clause, and \textit{begin} is the index of the most recently propagated 
literal in the resolution clause. So when the most recently and least recently propagated 
literals are the same, there must be only one left, therefore we can stop. Although it's 
clever, it's not efficient, and therefore, not very clever.

\subsection{First UIPs - Results}

Unfortunately our runtime increases again. At this point, I'm willing to accept
that perhaps I was just really really good at implementing DPLL, but since I'm not 
good at implementing CDCL (I literally failed the first time), my implementation 
is not up to par. In the next section, I will identify possible improvements that 
can be made, and hopefully they'll make a \textit{positive} difference.

\noindent
uf20.91: \textbf{406ms}

\noindent
UUF50.218.1000: \textbf{14012ms}

\begin{FVerbatim}
On all 1000 equations in uf20-91:

First UIPs:                406ms (~0.41 seconds)
Backjumping:               380ms (~0.38 seconds).
Clause Learning:           425ms (~0.43 seconds).
Iterative DPLL (AbP):      328ms (~0.33 seconds).
Iterative DPLL (PbA):      324ms (~0.32 seconds).
Watched Literals:          349ms (~0.35 seconds).
Pure Literal Elimination:  1548ms (~1.55 seconds).
Unit Propagation:          868ms (~0.87 seconds).
Recursive Backtracking:    1,540,613ms (~25.68 minutes).
Brute Force:               1,770,851ms (~29.51 minutes).
Short-Circuited (Basic):   49,874ms (~50 seconds).
Parallel:                  10,589ms (~10.5 seconds).

On all 1000 equations in UUF50.218.1000:

First UIPs:               14,012ms (~14.01 seconds).
Backjumping:              11,291ms (~11.29 seconds).
Clause Learning:          9,652ms (~9.65 seconds).
Iterative DPLL (AbP):     6,369ms (~6.37 seconds).
Iterative DPLL (PbA):     6,365ms (~6.37 seconds).
Watched Literals:         7,130ms (~7.13 seconds).
Pure Literal Elimination: 231,107ms (~3.85 minutes).
Unit Propagation:         140,041ms (~2.34 minutes).
\end{FVerbatim}